\section{Limitations and Future Work}
Our work significantly advances automated evaluation of LLM-based CTF agents, revealing insights into reasoning-retrieval interactions and improving partial scoring, solution quality, and cost-efficiency over baselines. Nevertheless, as this field is nascent, several opportunities exist for further refinement. The  CTF Competency Index employs a fixed, expert-curated weighting scheme that, while effective for studied challenges, could benefit from tuning or retraining to accommodate emerging or atypical CTF tasks. Our benchmarking focuses on a single open-source agent architecture (D-CIPHER) due to unavailability of alternative offensive-security frameworks; extending evaluations to fundamentally different designs would provide valuable insights. Moreover, dynamic metrics that adapt to domain drift would enhance the framework’s applicability to complex CTF scenarios. %These directions highlight promising avenues to extend this research’s generalizability and impact.

Several directions can strengthen our framework. First, adaptive calibration for the CCI—alongside automated online evaluation metrics that detect domain drift and adversarial shifts—would improve resilience and support dynamic reweighting for new challenge types. Second, expanding baseline agents to include reinforcement-learning agentic systems, multi-role division multi-agent setups, and knowledge-based retrieval will provide broader performance insights. Third, more rigorous targeted ablation studies isolating explicit chain-of-thought and implicit reasoning will clarify each approach’s contribution to success and robustness. %Pursuing these avenues will advance robust, versatile LLM-driven offensive security agents.